% =====================================================================
\section{Conclusión y Trabajo Futuro}\label{sec:conclusion}
% =====================================================================

Introdujimos \textbf{Walk Attention}, una arquitectura que mitiga el
over-squashing en redes neuronales de grafos atendiendo sobre un soporte
multi-salto del álgebra de caminos del quiver subyacente. El álgebra graduada
da, vía la Proposición~\ref{prop:count}, qué pares de nodos conectan caminos de
longitud-$g$ y con qué multiplicidad; con ello definimos una atención dispersa
de soporte algebraico y pesos aprendidos. Sobre un benchmark ajustable resolvió
exactamente una recuperación de largo alcance ($1.000\pm0.000$ sobre cinco
semillas) donde la atención de un salto colapsó al azar y un agregador
multi-salto de pesos fijos solo tuvo éxito parcial; en Peptides-func fue
competitivo con líneas base fuertes. Su ventaja depende del régimen: decisiva
bajo over-squashing severo y comparable a un salto cuando es leve. Un resultado
negativo mostró que comprimir caminos equivalentes con el \emph{cociente} es
contraproducente cuando la multiplicidad porta señal: el papel del álgebra es
determinar el soporte, no descartarla.

La construcción es elemental y autocontenida, como puente entre la teoría de
representaciones de quivers y el aprendizaje de representaciones de grafos. Walk
Attention es la lectura
aprendida y representacional de un sustrato algebraico común
(Sección~\ref{sec:ground})---el álgebra de caminos graduada y su sistema de
vecindades por grado de impacto---cuya lectura dinámica es el marco de autómatas de
impacto sobre quivers \cite{zainea2026aiq}; relacionarlas es ya una dirección, y
siguen varias más.

\begin{itemize}
  \item \textbf{Benchmarks con over-squashing severo.} Peptides-func es leve y el
  soporte multi-salto no es decisivo. El siguiente paso es buscar tareas reales de
  largo alcance con cuellos de botella agudos---donde la ventaja debería
  transferirse---y compararlo con rewiring \cite{topping2022understanding} y graph
  Transformers.
  \item \textbf{Soportes con forma de cociente.} El cociente $\kQ/I$, en lugar de
  comprimir la señal, puede \emph{particionar} los caminos en clases de
  equivalencia y enrutar clases distintas a cabezas distintas. Aprender qué $I$
  sirven---como el hermano dinámico las fija por regla---es atractivo.
  \item \textbf{Escalar el soporte.} En grafos densos el soporte de rango-$k$
  crece; truncarlo, muestrearlo o aprender uno más disperso acota el costo sin
  sacrificar largo alcance.
\end{itemize}

\subsubsection*{Reproducibilidad.}
Liberamos los generadores de grafos, los operadores de walk, la capa Walk
Attention, las líneas base y los registros de experimentos. Los cálculos del
álgebra usan una biblioteca abierta de quivers/álgebras de caminos; al ser su
primera aplicación publicada a GNNs, mantuvimos la construcción explícita.
